\makereporttitle{2024}{aprilie}

\makereportpreamble

\textbf{Activitate de cercetare:}

Plecând de la nevoia de a putea diagonaliza eficient o matrice simetrică de dimensiuni mari, care depinde de un parametru real, am început să parcurg și să înțeleg articolul „\emph{Calculating eigenvalues and eigenvectors of parameter-dependent hamiltonians using an adaptative wave operator method}” a lui Arnaud Leclerc și Georges Jolicard, publicat în 2020 în \emph{Journal of Chemical Physics}, volumul 152, numărul 20. Obiectivul meu este să adaptez metoda descrisă în lucrare pentru a diagonaliza matricea obținută discretizând Hamiltonianul cuadrupol bosonic, păstrând energia totală fixă și variind parametrul \(B\).

Algoritmul descris în lucrare folosește o serie de pași pentru a obține în mod iterativ valorile și vectorii proprii ale unei matrici, fiind apoi capabil să obțină rezultatele actualizate în momentul în care se variază parametrii de care depinde matricea. Pentru a înțelege mai bine de unde vine ideea din spatele metodei și termenii folosiți, am început să aprofundez noțiuni de algebră liniară numerică: descompuneri QR, proiectori, transformări Householder, rotații Givens etc.

\vspace{0.5em}

\textbf{Cursuri și seminare științifice:}

În această lună am participat la cursul doctoral „\emph{Metodologia cercetării științifice și prelucrarea statistică a datelor}”, ținut de conf.\ dr.\ Cristian Necula. În cadrul acestui curs ne-au fost prezentate noțiuni introductive de analiză a datelor și de prelucrarea seriilor de timp, precum și noțiuni mai avansate cum ar fi analiza semnalelor în domeniul frecvență folosind transformata Fourier sau transforma \emph{wavelet}.

De câte ori am avut ocazia, am fost și la cursul de Fizică Cuantică ținut de dl.\ prof.\ Virgil Băran la Facultatea de Matematică și Informatică. Am consolidat în acest fel cunoștințele mele despre postulatele mecanicii cuantice și formalismul matematic asociat.

Am participat la câteva seminare științifice care s-au desfășurat la Institutul de Matematică al Academiei Române (IMAR):
\begin{itemize}
    \item Seminarul de geometrie organizat de Sergiu Moroianu, la o prelegere despre „Sisteme integrabile pe varietăți Poisson. Ecuația pereche Lax”;
    \item Seminarul ,,\emph{Analysis \& PDEs}'', organizat de Liviu Ignat și Cristian Cazacu, la o prelegere despre rezolvarea ecuației căldurii și comportamentul asimptotic al acesteia.
\end{itemize}

\vspace{0.5em}

\textbf{Observațiile conducătorului științific}:

\vspace{3em}

\makereportpostamble{29 aprilie 2024}
